\section{Limitations and Scope}

The theory makes one approximation explicit: the broadband surrogate remains the only approximation step in the chain. The scaling law \(\tau^*(L)=d_{\mathrm{head}}/\sqrt{L}\) is therefore treated as an empirical law / conjecture rather than as an unconditional theorem. This separation is deliberate. The exact claim concerns the surrogate-to-ODE derivation; the predictive claim concerns how the effective optimum evolves with training length.

On the empirical side, the strongest primary downstream evidence is still retrieval-heavy. The \(750\mathrm{M}\) continue experiment is useful because it shows the same short-range-cost / long-range-gain pattern at larger scale, but it remains single-seed supporting evidence, and its long-form downstream comparison was interrupted by a dataset-download failure. The video temporal results are encouraging and high-upside, but they remain supportive rather than co-primary at the current evidence level.

We also make a deliberate methodological choice about what \emph{not} to treat as primary evidence. Short-step LoRA or instruction-tuning results can be suggestive, but they confound positional adaptation with task adaptation, which makes them weak tests of whether the frequency allocation itself is responsible for the gain. For the same reason, not every positive larger-scale or cross-modal signal is elevated to headline status. Multi-seed text results remain the primary anchors, while single-seed scaling runs and video transfer are used to test whether the same qualitative mechanism continues to hold outside the core setting.
